\section{ANALYTICAL SOLUTIONS}
\subsection{NORMALISED SYSTEM}

Aligning with the goals of this research we will show the analytical solutions for a singular problem setup across
five integration schemes under ten timesteps.

Before starting any calculations, it is required to apply the appropriate units of measurement.
It was decided to apply a normalised system due to the ease of debugging and interpretation, numerical stability
and avoiding the vulnerability of floating-point errors.

Hence why, the units are as follows:
\begin{itemize}
    \item Distance -- Astronomical Units (Earth to Sun Distance)
    \item Mass -- Solar Masses
    \item Gravitational Attraction -- $4\pi^2$
    \item Timestep -- Years
\end{itemize}
Following studies a timestep of $\Delta t = 10^{-3}$\,yr has been selected, since unstable integrators tend to
converge later in the simulation producing more reasonable data.
This allows tracking of convergences and energy preservation more accurately.

Due to the assumption of a normalised numerical system the value of Gravitational Attraction Constant ($G$) has to be
derived from those systemic values.
Physics remains the same, while the measuring system modifies.
Since the normalised system is based on the relationship between the Earth and the Sun, we can derive the value of $G$
that complies with the numerical system.

Applying Newtonian gravity,

\begin{equation}
    F = G\frac{m_1 m_2}{r^2},
\end{equation}

and the equation for circular motion,

\begin{equation}
    F = \frac{m_2 v^2}{r},
\end{equation}

we obtain

\begin{equation}
    \frac{m_2 v^2}{r}
    =
    G\frac{m_1 m_2}{r^2}.
\end{equation}

Cancelling $m_2$ from both sides yields

\begin{equation}
    v^2 = \frac{Gm_1}{r}.
\end{equation}

Since the circumference of a circular orbit is $2\pi r$, the orbital velocity can be expressed in terms of the orbital period $T$ as

\begin{equation}
    v = \frac{2\pi r}{T}.
\end{equation}

Substituting this expression for $v$ into Equation~(3) gives

\begin{equation}
    \left(\frac{2\pi r}{T}\right)^2 = \frac{Gm_1}{r}.
\end{equation}

Solving for $T$ yields

\begin{equation}
    T^2 = \frac{4\pi^2 r^3}{Gm_1}.
\end{equation}

Assuming a normalised system of units in which the orbital period, orbital radius, and central mass are all equal to
unity ($T = r = m_1 = 1$), we obtain

\begin{equation}
    1 = \frac{4\pi^2}{G}.
\end{equation}

Rearranging gives

\begin{equation}
    G = 4\pi^2.
\end{equation}

\subsection{ANALYTICAL PROBLEM}

The problem definition is as follows: A star of one solar mass $M = 1\,M_\odot$ is fixed at the origin.
A planet orbiting the star is $333000^{-1}$ of solar mass and is placed at initial position
$\vec{r}_0 = (2,\,1)^T$\,AU.
All quantities are expressed in natural astronomical units: distances in AU, time in years, and mass in
solar masses, with $G = 4\pi^2\,\mathrm{AU}^3\,\mathrm{yr}^{-2}\,M_\odot^{-1}$.
The simulation procedurally generates two objects containing their mass and coordinates, while the rest of the
calculations are presented through the later simulation.

\subsection{INITIAL VALUES}

\subsubsection*{Initial distance}

The initial separation between the planet and the star is

\begin{equation}
    r_0 = |\vec{r}_0| = \sqrt{x_0^2 + y_0^2} = \sqrt{2^2 + 1^2} = \sqrt{5} \approx 2.23606798\,\mathrm{AU}.
\end{equation}

\subsubsection*{Initial acceleration}

Newton's law of gravity gives the acceleration per unit planet mass as

\begin{equation}
    \vec{a} = -\frac{GM}{r^3}\vec{r}.
\end{equation}

Substituting $G = 4\pi^2$, $M = 1$, $r_0 = \sqrt{5}$, $\vec{r}_0 = (2,1)^T$:

\begin{equation}
    \vec{a}_0 = -\frac{4\pi^2}{(\sqrt{5})^3}\begin{pmatrix}2\\1\end{pmatrix}
    = -\frac{4\pi^2}{5\sqrt{5}}\begin{pmatrix}2\\1\end{pmatrix}
    \approx \begin{pmatrix}-7.06211403\\-3.53105702\end{pmatrix}\,\mathrm{AU\,yr}^{-2}.
\end{equation}

The magnitude is $|\vec{a}_0| = 4\pi^2/5 \approx 7.8957\,\mathrm{AU\,yr}^{-2}$.

\subsubsection*{Circular orbital velocity}

For a circular orbit, gravitational acceleration exactly supplies centripetal acceleration:

\begin{equation}
    v_c = \sqrt{\frac{GM}{r_0}} = \sqrt{\frac{4\pi^2}{\sqrt{5}}} \approx 4.20181926\,\mathrm{AU\,yr}^{-1}.
\end{equation}

The velocity is directed tangentially (perpendicular to $\vec{r}_0$, counter-clockwise):

\begin{equation}
    \hat{t} = \frac{(-y_0,\,x_0)^T}{r_0} = \frac{1}{\sqrt{5}}\begin{pmatrix}-1\\2\end{pmatrix},
    \qquad
    \vec{v}_0 = v_c\,\hat{t} \approx \begin{pmatrix}-1.87911070\\3.75822140\end{pmatrix}\,\mathrm{AU\,yr}^{-1}.
\end{equation}

\subsubsection*{Initial conserved quantities}

The total mechanical energy (kinetic plus potential) per unit planet mass is

\begin{equation}
    E_0 = \frac{1}{2}m|\vec{v}_0|^2 - \frac{GMm}{r_0}
    = -\frac{GMm}{2r_0}
    \approx -2.6509\times10^{-5}\,M_\odot\,\mathrm{AU}^2\,\mathrm{yr}^{-2},
\end{equation}

and the scalar $z$-component of angular momentum is

\begin{equation}
    L_0 = m\!\left(x_0 v_{y,0} - y_0 v_{x,0}\right)
    \approx 2.8215\times10^{-5}\,M_\odot\,\mathrm{AU}^2\,\mathrm{yr}^{-1}.
\end{equation}

\subsection{FIRST TIMESTEP DERIVATION ($\Delta t = 10^{-3}$\,yr)}

We adopt a time step $\Delta t = 10^{-3}$\,yr and derive the first update for each integrator explicitly.
Relative errors in the conserved quantities are defined as
$\varepsilon_E = |E_n - E_0|/|E_0|$ and $\varepsilon_L = |L_n - L_0|/|L_0|$.

\subsubsection*{Explicit Euler}

Update rule: $\vec{r}_{n+1} = \vec{r}_n + \Delta t\,\vec{v}_n$,\quad $\vec{v}_{n+1} = \vec{v}_n + \Delta t\,\vec{a}_n$.

Step $n = 0 \to 1$:
\begin{align*}
    \vec{r}_1 &= \begin{pmatrix}2\\1\end{pmatrix}
    + 10^{-3}\begin{pmatrix}-1.87911070\\3.75822140\end{pmatrix}
    = \begin{pmatrix}1.99812089\\1.00375822\end{pmatrix}\,\mathrm{AU},\\
    \vec{v}_1 &= \begin{pmatrix}-1.87911070\\3.75822140\end{pmatrix}
    + 10^{-3}\begin{pmatrix}-7.06211403\\-3.53105702\end{pmatrix}
    = \begin{pmatrix}-1.88617281\\3.75469034\end{pmatrix}\,\mathrm{AU\,yr}^{-1}.
\end{align*}
Both position and velocity are updated using values from the same timestep, so this method is first-order
accurate and does not conserve energy over time.

\subsubsection*{Symplectic Euler}

Update rule: velocity is updated first, then position uses the updated velocity:
$\vec{v}_{n+1} = \vec{v}_n + \Delta t\,\vec{a}_n$,\quad $\vec{r}_{n+1} = \vec{r}_n + \Delta t\,\vec{v}_{n+1}$.

Step $n = 0 \to 1$:
\begin{align*}
    \vec{v}_1 &= \begin{pmatrix}-1.87911070\\3.75822140\end{pmatrix}
    + 10^{-3}\begin{pmatrix}-7.06211403\\-3.53105702\end{pmatrix}
    = \begin{pmatrix}-1.88617281\\3.75469034\end{pmatrix}\,\mathrm{AU\,yr}^{-1},\\
    \vec{r}_1 &= \begin{pmatrix}2\\1\end{pmatrix}
    + 10^{-3}\begin{pmatrix}-1.88617281\\3.75469034\end{pmatrix}
    = \begin{pmatrix}1.99811383\\1.00375469\end{pmatrix}\,\mathrm{AU}.
\end{align*}
Symplectic Euler is also first-order, but it preserves the symplectic structure of Hamiltonian mechanics, resulting
in bounded (rather than drifting) energy errors.

\subsubsection*{Velocity Verlet}

Update rule: $\vec{r}_{n+1} = \vec{r}_n + \Delta t\,\vec{v}_n + \tfrac{1}{2}\Delta t^2\,\vec{a}_n$,\quad
$\vec{v}_{n+1} = \vec{v}_n + \tfrac{\Delta t}{2}(\vec{a}_n + \vec{a}_{n+1})$.

Step $n = 0 \to 1$:
\begin{align*}
    \vec{r}_1 &= \begin{pmatrix}2\\1\end{pmatrix}
    + 10^{-3}\begin{pmatrix}-1.87911070\\3.75822140\end{pmatrix}
    + \frac{10^{-6}}{2}\begin{pmatrix}-7.06211403\\-3.53105702\end{pmatrix}
    = \begin{pmatrix}1.99811736\\1.00375646\end{pmatrix}\,\mathrm{AU}.
\end{align*}
Then $\vec{a}_1$ is recomputed from $\vec{r}_1$, and
\begin{equation*}
    \vec{v}_1 = \begin{pmatrix}-1.87911070\\3.75822140\end{pmatrix}
    + \frac{10^{-3}}{2}
    \left[\begin{pmatrix}-7.06211403\\-3.53105702\end{pmatrix}
        + \vec{a}_1\right]
    = \begin{pmatrix}-1.88616949\\3.75468371\end{pmatrix}\,\mathrm{AU\,yr}^{-1}.
\end{equation*}

\subsubsection*{Leapfrog (Kick--Drift--Kick)}

The Leapfrog integrator stores velocities at half-integer steps.
The initialisation half-kick is:
\begin{equation*}
    \vec{v}_{1/2} = \vec{v}_0 + \frac{\Delta t}{2}\vec{a}_0
    = \begin{pmatrix}-1.87911070\\3.75822140\end{pmatrix}
    + 5\times10^{-4}\begin{pmatrix}-7.06211403\\-3.53105702\end{pmatrix}
    = \begin{pmatrix}-1.88264677\\3.75645587\end{pmatrix}\,\mathrm{AU\,yr}^{-1}.
\end{equation*}
Full drift step $n = 0 \to 1$:
\begin{equation*}
    \vec{r}_1 = \vec{r}_0 + \Delta t\,\vec{v}_{1/2}
    = \begin{pmatrix}2\\1\end{pmatrix}
    + 10^{-3}\begin{pmatrix}-1.88264677\\3.75645587\end{pmatrix}
    = \begin{pmatrix}1.99811736\\1.00375646\end{pmatrix}\,\mathrm{AU}.
\end{equation*}
Velocities at integer steps are recovered as $\vec{v}_n = \vec{v}_{n+1/2} - \tfrac{\Delta t}{2}\vec{a}_n$,
producing results identical to Velocity Verlet to floating-point precision.

\subsubsection*{RK4 (4th-order Runge--Kutta)}

Writing the state vector as $\mathbf{y} = (\vec{r},\,\vec{v})^T$ with
$\dot{\mathbf{y}} = (\vec{v},\,\vec{a}(\vec{r}))^T$, four intermediate slopes are:
\begin{equation*}
    \mathbf{k}_1 = f(\mathbf{y}_n),\quad
    \mathbf{k}_2 = f\!\left(\mathbf{y}_n + \tfrac{\Delta t}{2}\mathbf{k}_1\right),\quad
    \mathbf{k}_3 = f\!\left(\mathbf{y}_n + \tfrac{\Delta t}{2}\mathbf{k}_2\right),\quad
    \mathbf{k}_4 = f(\mathbf{y}_n + \Delta t\,\mathbf{k}_3),
\end{equation*}
\begin{equation*}
    \mathbf{y}_{n+1} = \mathbf{y}_n + \frac{\Delta t}{6}\!\left(\mathbf{k}_1 + 2\mathbf{k}_2 + 2\mathbf{k}_3 + \mathbf{k}_4\right).
\end{equation*}
After step $n = 0 \to 1$:
\begin{equation*}
    \vec{r}_1 \approx \begin{pmatrix}1.99811736\\1.00375645\end{pmatrix}\,\mathrm{AU},\qquad
    \vec{v}_1 \approx \begin{pmatrix}-1.88616949\\3.75468371\end{pmatrix}\,\mathrm{AU\,yr}^{-1}.
\end{equation*}
RK4 is fourth-order accurate; energy errors at each step are $\mathcal{O}(\Delta t^5)$.

\subsection{10-STEP SIMULATION TABLES ($\Delta t = 10^{-3}$\,yr)}

The conserved quantities and their relative errors are defined as
\begin{equation}
    E_n = \frac{1}{2}m\!\left(v_{x,n}^2 + v_{y,n}^2\right) - \frac{GMm}{r_n},
    \quad
    L_n = m\!\left(x_n v_{y,n} - y_n v_{x,n}\right),
    \quad
    \varepsilon_E = \frac{|E_n - E_0|}{|E_0|},
    \quad
    \varepsilon_L = \frac{|L_n - L_0|}{|L_0|}.
\end{equation}
Units: positions and distances in AU, velocities in AU\,yr$^{-1}$, accelerations in AU\,yr$^{-2}$,
energies in $M_\odot\,\mathrm{AU}^2\,\mathrm{yr}^{-2}$, angular momenta in $M_\odot\,\mathrm{AU}^2\,\mathrm{yr}^{-1}$.

\subsubsection*{Explicit Euler}

\begin{table}[H]
    \centering
    \scriptsize
    \setlength{\tabcolsep}{4pt}
    \begin{tabular}{c l l l r r r r}
        \hline
        $n$ & $\vec{r}$ (AU) & $\vec{v}$ (AU\,yr$^{-1}$) & $\vec{a}$ (AU\,yr$^{-2}$) & $|\vec{v}|$ & $r$ & $\varepsilon_E$ & $\varepsilon_L$ \\
        \hline
        0  & (2.00, 1.00) & ($-$1.88, 3.76)  & ($-$7.06, $-$3.53) & 4.20 & 2.24 & $0$ & $0$ \\
        1  & (2.00, 1.00) & ($-$1.89, 3.75)  & ($-$7.06, $-$3.54) & 4.20 & 2.24 & $7.06\!\times\!10^{-6}$  & $3.53\!\times\!10^{-6}$ \\
        2  & (2.00, 1.01) & ($-$1.89, 3.75)  & ($-$7.05, $-$3.56) & 4.20 & 2.24 & $1.41\!\times\!10^{-5}$  & $7.06\!\times\!10^{-6}$ \\
        3  & (1.99, 1.01) & ($-$1.90, 3.75)  & ($-$7.04, $-$3.57) & 4.20 & 2.24 & $2.12\!\times\!10^{-5}$  & $1.06\!\times\!10^{-5}$ \\
        4  & (1.99, 1.02) & ($-$1.91, 3.74)  & ($-$7.04, $-$3.58) & 4.20 & 2.24 & $2.82\!\times\!10^{-5}$  & $1.41\!\times\!10^{-5}$ \\
        5  & (1.99, 1.02) & ($-$1.91, 3.74)  & ($-$7.03, $-$3.60) & 4.20 & 2.24 & $3.53\!\times\!10^{-5}$  & $1.77\!\times\!10^{-5}$ \\
        6  & (1.99, 1.02) & ($-$1.92, 3.74)  & ($-$7.02, $-$3.61) & 4.20 & 2.24 & $4.24\!\times\!10^{-5}$  & $2.12\!\times\!10^{-5}$ \\
        7  & (1.99, 1.03) & ($-$1.93, 3.73)  & ($-$7.01, $-$3.62) & 4.20 & 2.24 & $4.94\!\times\!10^{-5}$  & $2.47\!\times\!10^{-5}$ \\
        8  & (1.98, 1.03) & ($-$1.94, 3.73)  & ($-$7.01, $-$3.64) & 4.20 & 2.24 & $5.65\!\times\!10^{-5}$  & $2.82\!\times\!10^{-5}$ \\
        9  & (1.98, 1.03) & ($-$1.94, 3.73)  & ($-$7.00, $-$3.65) & 4.20 & 2.24 & $6.36\!\times\!10^{-5}$  & $3.18\!\times\!10^{-5}$ \\
        10 & (1.98, 1.04) & ($-$1.95, 3.72)  & ($-$6.99, $-$3.66) & 4.20 & 2.24 & $7.06\!\times\!10^{-5}$  & $3.53\!\times\!10^{-5}$ \\
        \hline
    \end{tabular}
    \caption{Explicit Euler, $\Delta t = 10^{-3}$\,yr. Energy and angular momentum drift monotonically at $\mathcal{O}(\Delta t)$ per step -- characteristic of a non-symplectic first-order method.}
\end{table}

\subsubsection*{Symplectic Euler}

\begin{table}[H]
    \centering
    \scriptsize
    \setlength{\tabcolsep}{4pt}
    \begin{tabular}{c l l l r r r r}
        \hline
        $n$ & $\vec{r}$ (AU) & $\vec{v}$ (AU\,yr$^{-1}$) & $\vec{a}$ (AU\,yr$^{-2}$) & $|\vec{v}|$ & $r$ & $\varepsilon_E$ & $\varepsilon_L$ \\
        \hline
        0  & (2.00, 1.00) & ($-$1.88, 3.76)  & ($-$7.06, $-$3.53) & 4.20 & 2.24 & $0$ & $0$ \\
        1  & (2.00, 1.00) & ($-$1.89, 3.75)  & ($-$7.06, $-$3.54) & 4.20 & 2.24 & $3.10\!\times\!10^{-12}$ & $\sim\!0$ \\
        2  & (2.00, 1.01) & ($-$1.89, 3.75)  & ($-$7.05, $-$3.56) & 4.20 & 2.24 & $1.30\!\times\!10^{-11}$ & $\sim\!0$ \\
        3  & (1.99, 1.01) & ($-$1.90, 3.75)  & ($-$7.04, $-$3.57) & 4.20 & 2.24 & $2.80\!\times\!10^{-11}$ & $\sim\!0$ \\
        4  & (1.99, 1.01) & ($-$1.91, 3.74)  & ($-$7.04, $-$3.58) & 4.20 & 2.24 & $5.00\!\times\!10^{-11}$ & $\sim\!0$ \\
        5  & (1.99, 1.02) & ($-$1.91, 3.74)  & ($-$7.03, $-$3.60) & 4.20 & 2.24 & $7.80\!\times\!10^{-11}$ & $\sim\!0$ \\
        6  & (1.99, 1.02) & ($-$1.92, 3.74)  & ($-$7.02, $-$3.61) & 4.20 & 2.24 & $1.10\!\times\!10^{-10}$ & $\sim\!0$ \\
        7  & (1.99, 1.03) & ($-$1.93, 3.73)  & ($-$7.02, $-$3.62) & 4.20 & 2.24 & $1.50\!\times\!10^{-10}$ & $\sim\!0$ \\
        8  & (1.98, 1.03) & ($-$1.94, 3.73)  & ($-$7.01, $-$3.64) & 4.20 & 2.24 & $2.00\!\times\!10^{-10}$ & $\sim\!0$ \\
        9  & (1.98, 1.03) & ($-$1.94, 3.73)  & ($-$7.00, $-$3.65) & 4.20 & 2.24 & $2.50\!\times\!10^{-10}$ & $\sim\!0$ \\
        10 & (1.98, 1.04) & ($-$1.95, 3.72)  & ($-$6.99, $-$3.66) & 4.20 & 2.24 & $3.10\!\times\!10^{-10}$ & $\sim\!0$ \\
        \hline
    \end{tabular}
    \caption{Symplectic Euler, $\Delta t = 10^{-3}$\,yr. Angular momentum is conserved to machine precision; energy error is bounded ($\sim\!10^{-10}$--$10^{-12}$) with no secular drift -- hallmark of a symplectic method.}
\end{table}

\subsubsection*{Velocity Verlet}

\begin{table}[H]
    \centering
    \scriptsize
    \setlength{\tabcolsep}{4pt}
    \begin{tabular}{c l l l r r r r}
        \hline
        $n$ & $\vec{r}$ (AU) & $\vec{v}$ (AU\,yr$^{-1}$) & $\vec{a}$ (AU\,yr$^{-2}$) & $|\vec{v}|$ & $r$ & $\varepsilon_E$ & $\varepsilon_L$ \\
        \hline
        0  & (2.00, 1.00) & ($-$1.88, 3.76)  & ($-$7.06, $-$3.53) & 4.20 & 2.24 & $0$ & $0$ \\
        1  & (2.00, 1.00) & ($-$1.89, 3.75)  & ($-$7.06, $-$3.54) & 4.20 & 2.24 & $\sim\!10^{-16}$ & $\sim\!0$ \\
        2  & (2.00, 1.01) & ($-$1.89, 3.75)  & ($-$7.05, $-$3.56) & 4.20 & 2.24 & $\sim\!10^{-16}$ & $\sim\!0$ \\
        3  & (1.99, 1.01) & ($-$1.90, 3.75)  & ($-$7.04, $-$3.57) & 4.20 & 2.24 & $\sim\!10^{-16}$ & $\sim\!0$ \\
        4  & (1.99, 1.02) & ($-$1.91, 3.74)  & ($-$7.04, $-$3.58) & 4.20 & 2.24 & $\sim\!10^{-16}$ & $\sim\!0$ \\
        5  & (1.99, 1.02) & ($-$1.91, 3.74)  & ($-$7.03, $-$3.60) & 4.20 & 2.24 & $\sim\!10^{-16}$ & $\sim\!0$ \\
        6  & (1.99, 1.02) & ($-$1.92, 3.74)  & ($-$7.02, $-$3.61) & 4.20 & 2.24 & $\sim\!10^{-16}$ & $\sim\!0$ \\
        7  & (1.99, 1.03) & ($-$1.93, 3.73)  & ($-$7.02, $-$3.62) & 4.20 & 2.24 & $\sim\!10^{-16}$ & $\sim\!0$ \\
        8  & (1.98, 1.03) & ($-$1.94, 3.73)  & ($-$7.01, $-$3.64) & 4.20 & 2.24 & $\sim\!10^{-16}$ & $\sim\!0$ \\
        9  & (1.98, 1.03) & ($-$1.94, 3.73)  & ($-$7.00, $-$3.65) & 4.20 & 2.24 & $\sim\!10^{-16}$ & $\sim\!0$ \\
        10 & (1.98, 1.04) & ($-$1.95, 3.72)  & ($-$6.99, $-$3.66) & 4.20 & 2.24 & $1.00\!\times\!10^{-15}$ & $\sim\!0$ \\
        \hline
    \end{tabular}
    \caption{Velocity Verlet, $\Delta t = 10^{-3}$\,yr. Both $|\vec{v}|$ and $r$ remain constant to 3 significant figures throughout; energy errors are at floating-point rounding level ($\sim\!10^{-15}$--$10^{-16}$). This is a second-order symplectic method.}
\end{table}

\subsubsection*{Leapfrog (Kick--Drift--Kick)}

\begin{table}[H]
    \centering
    \scriptsize
    \setlength{\tabcolsep}{4pt}
    \begin{tabular}{c l l l r r r r}
        \hline
        $n$ & $\vec{r}$ (AU) & $\vec{v}$ (AU\,yr$^{-1}$) & $\vec{a}$ (AU\,yr$^{-2}$) & $|\vec{v}|$ & $r$ & $\varepsilon_E$ & $\varepsilon_L$ \\
        \hline
        0  & (2.00, 1.00) & ($-$1.88, 3.76)  & ($-$7.06, $-$3.53) & 4.20 & 2.24 & $0$ & $0$ \\
        1  & (2.00, 1.00) & ($-$1.89, 3.75)  & ($-$7.06, $-$3.54) & 4.20 & 2.24 & $\sim\!0$    & $\sim\!0$ \\
        2  & (2.00, 1.01) & ($-$1.89, 3.75)  & ($-$7.05, $-$3.56) & 4.20 & 2.24 & $\sim\!10^{-16}$ & $\sim\!0$ \\
        3  & (1.99, 1.01) & ($-$1.90, 3.75)  & ($-$7.04, $-$3.57) & 4.20 & 2.24 & $\sim\!10^{-16}$ & $\sim\!0$ \\
        4  & (1.99, 1.02) & ($-$1.91, 3.74)  & ($-$7.04, $-$3.58) & 4.20 & 2.24 & $\sim\!10^{-16}$ & $\sim\!0$ \\
        5  & (1.99, 1.02) & ($-$1.91, 3.74)  & ($-$7.03, $-$3.60) & 4.20 & 2.24 & $\sim\!10^{-16}$ & $\sim\!0$ \\
        6  & (1.99, 1.02) & ($-$1.92, 3.74)  & ($-$7.02, $-$3.61) & 4.20 & 2.24 & $\sim\!10^{-16}$ & $\sim\!0$ \\
        7  & (1.99, 1.03) & ($-$1.93, 3.73)  & ($-$7.02, $-$3.62) & 4.20 & 2.24 & $\sim\!10^{-16}$ & $\sim\!0$ \\
        8  & (1.98, 1.03) & ($-$1.94, 3.73)  & ($-$7.01, $-$3.64) & 4.20 & 2.24 & $\sim\!10^{-16}$ & $\sim\!0$ \\
        9  & (1.98, 1.03) & ($-$1.94, 3.73)  & ($-$7.00, $-$3.65) & 4.20 & 2.24 & $\sim\!10^{-16}$ & $\sim\!0$ \\
        10 & (1.98, 1.04) & ($-$1.95, 3.72)  & ($-$6.99, $-$3.66) & 4.20 & 2.24 & $\sim\!10^{-16}$ & $\sim\!0$ \\
        \hline
    \end{tabular}
    \caption{Leapfrog (KDK), $\Delta t = 10^{-3}$\,yr. Results are numerically identical to Velocity Verlet (they are algebraically equivalent for constant-mass conservative systems). The Leapfrog is distinguished by its staggered storage of velocities at half-integer steps.}
\end{table}

\subsubsection*{RK4 (4th-order Runge--Kutta)}

\begin{table}[H]
    \centering
    \scriptsize
    \setlength{\tabcolsep}{4pt}
    \begin{tabular}{c l l l r r r r}
        \hline
        $n$ & $\vec{r}$ (AU) & $\vec{v}$ (AU\,yr$^{-1}$) & $\vec{a}$ (AU\,yr$^{-2}$) & $|\vec{v}|$ & $r$ & $\varepsilon_E$ & $\varepsilon_L$ \\
        \hline
        0  & (2.00, 1.00) & ($-$1.88, 3.76)  & ($-$7.06, $-$3.53) & 4.20 & 2.24 & $0$ & $0$ \\
        1  & (2.00, 1.00) & ($-$1.89, 3.75)  & ($-$7.06, $-$3.54) & 4.20 & 2.24 & $\sim\!0$ & $\sim\!0$ \\
        2  & (2.00, 1.01) & ($-$1.89, 3.75)  & ($-$7.05, $-$3.56) & 4.20 & 2.24 & $\sim\!0$ & $\sim\!0$ \\
        3  & (1.99, 1.01) & ($-$1.90, 3.75)  & ($-$7.04, $-$3.57) & 4.20 & 2.24 & $\sim\!0$ & $\sim\!0$ \\
        4  & (1.99, 1.02) & ($-$1.91, 3.74)  & ($-$7.04, $-$3.58) & 4.20 & 2.24 & $\sim\!0$ & $\sim\!0$ \\
        5  & (1.99, 1.02) & ($-$1.91, 3.74)  & ($-$7.03, $-$3.60) & 4.20 & 2.24 & $\sim\!0$ & $\sim\!0$ \\
        6  & (1.99, 1.02) & ($-$1.92, 3.74)  & ($-$7.02, $-$3.61) & 4.20 & 2.24 & $\sim\!10^{-16}$ & $\sim\!0$ \\
        7  & (1.99, 1.03) & ($-$1.93, 3.73)  & ($-$7.02, $-$3.62) & 4.20 & 2.24 & $\sim\!10^{-16}$ & $\sim\!0$ \\
        8  & (1.98, 1.03) & ($-$1.94, 3.73)  & ($-$7.01, $-$3.64) & 4.20 & 2.24 & $\sim\!0$ & $\sim\!0$ \\
        9  & (1.98, 1.03) & ($-$1.94, 3.73)  & ($-$7.00, $-$3.65) & 4.20 & 2.24 & $\sim\!0$ & $\sim\!0$ \\
        10 & (1.98, 1.04) & ($-$1.95, 3.72)  & ($-$6.99, $-$3.66) & 4.20 & 2.24 & $\sim\!0$ & $\sim\!0$ \\
        \hline
    \end{tabular}
    \caption{RK4, $\Delta t = 10^{-3}$\,yr. Both $|\vec{v}|$ and $r$ are constant to 3 significant figures throughout; energy errors are entirely at floating-point rounding level. RK4 is not symplectic but is 4th-order accurate, so local errors are $\mathcal{O}(\Delta t^5)$.}
\end{table}

\subsection{SUMMARY OF INTEGRATOR PERFORMANCE}

\begin{table}[H]
    \centering
    \begin{tabular}{l c c l}
        \hline
        Integrator & Order & $\varepsilon_E$ at $n=10$ & Conservation behaviour \\
        \hline
        Explicit Euler   & 1 & $7.06\times10^{-5}$ & Monotonic drift (non-symplectic) \\
        Symplectic Euler & 1 & $3.10\times10^{-10}$ & Bounded oscillation (symplectic) \\
        Velocity Verlet  & 2 & $\sim10^{-15}$      & Near machine precision (symplectic) \\
        Leapfrog         & 2 & $\sim10^{-16}$      & Identical to Velocity Verlet \\
        RK4              & 4 & $\sim10^{-16}$      & Near machine precision (not symplectic) \\
        \hline
    \end{tabular}
    \caption{Energy relative error $\varepsilon_E$ at $n = 10$ for each integrator with $\Delta t = 10^{-3}$\,yr.}
\end{table}

Key observations:

\begin{itemize}
    \item \textbf{Explicit Euler} loses energy monotonically at rate $\mathcal{O}(\Delta t)$ per step -- unacceptable for long-term orbital integration. At the finer $\Delta t = 10^{-3}$, however, the error at $n=10$ is $7.06\times10^{-5}$, two orders of magnitude smaller than at $\Delta t = 10^{-2}$.
    \item \textbf{Symplectic Euler} oscillates around the true energy with no secular drift, despite having the same first-order cost. Energy error ($\sim\!3\times10^{-10}$) is orders of magnitude lower than Explicit Euler; angular momentum is conserved to machine precision.
    \item \textbf{Velocity Verlet} and \textbf{Leapfrog} are algebraically equivalent and provide near-machine-precision conservation at second-order cost -- the standard choice in $N$-body simulation. At $\Delta t = 10^{-3}$ the energy error is entirely at the floating-point rounding floor ($\sim\!10^{-15}$--$10^{-16}$).
    \item \textbf{RK4} gives the lowest error per step at this timestep, with energy errors indistinguishable from machine epsilon. However, RK4 is not symplectic, so secular energy drift will appear on very long timescales; it also requires four force evaluations per step (versus one for Euler and two for Verlet).
\end{itemize}